\RequirePackage{pdfmanagement}
\SetKeys[document/metadata]{pdfversion=1.7, pdfstandard=A-2u, lang=en-GB}
\input{preamble/pdfA-compliance/glyphtounicode}
\documentclass[11pt,a4paper]{report}

\newcommand*{\PaperTitle}{No Andronov--Hopf bifurcation, no periodic orbit: cases~9 and~10 of Banaji--Boros--Hofbauer}
\newcommand*{\PaperTitlePlain}{No Andronov-Hopf bifurcation, no periodic orbit: cases 9 and 10 of Banaji-Boros-Hofbauer}
\newcommand*{\PaperShortTitle}{BBH cases 9 and 10}
\newcommand*{\PaperSubtitle}{chemistry-open-problems \textperiodcentered{} 24 October 2026 \textperiodcentered{} preliminary (draft by the solver, not compiled)}
\newcommand*{\PaperCredits}{Solver: Claude Opus 5.5 (effort xhigh) \textperiodcentered{} Independent verification: pending}
\newcommand*{\PaperKeywords}{mass-action kinetics, limit cycles, Li\'enard systems, Andronov--Hopf bifurcation, quadratic systems}
\input{chemop-preamble}
\addbibresource{references.bib}

% \R (= \mathbb{R}) is provided by the shared preamble
\newcommand{\Divop}{\operatorname{div}}
\newcommand{\Trop}{\operatorname{tr}}
\newcommand{\Gt}{\widetilde G}

\begin{document}
\makepapertitle

\begin{paperabstract}
Banaji, Boros and Hofbauer classified the two-species, three-reaction mass-action networks with at most bimolecular
sources that admit an Andronov--Hopf bifurcation. The two configurations that do are cases~9 and~10, with source
complexes \(2X, X+Y, Y\) and \(2X, X+Y, 0\) and arbitrary targets. They asked whether the networks of these cases that
admit no Andronov--Hopf bifurcation can nevertheless have a limit cycle. We show that they cannot: for every such
network and every choice of rate constants, the mass-action system has no periodic orbit in the closed positive
quadrant. Hence, granting their supercriticality computation, a network of case~9 or~10 has a limit cycle for some
rate constants if and only if it admits an Andronov--Hopf bifurcation. The proof first reduces everything to three
explicit families. These are linear in~\(y\), and become Li\'enard systems after a shear along the \(x\)-nullcline. A
comparison lemma for Li\'enard systems then reduces non-existence to a single inequality that does not involve the
rate constants, which we prove in closed form.
\end{paperabstract}

\section*{The question}

Banaji, Boros and Hofbauer~\autocite{BBH2024} (BBH) list the ten configurations of three non-collinear, at most
bimolecular source complexes on two species. Their Theorem~9 (arXiv numbering; Theorem~3 of the published version,
with identical text) decides, for each configuration, which networks admit an Andronov--Hopf bifurcation. Only
cases~9 (\(2X, X+Y, Y\)) and~10 (\(2X, X+Y, 0\)) contain such networks. In their Conclusions (§6) they write:

\begin{openquote}
``However, we leave it open whether the networks in cases 9 and 10 admitting no Andronov--Hopf bifurcations can have
a limit cycle. For example, we do not know whether the system \(2X\to3X+2Y\), \(X+Y\to0\), \(Y\to X+Y\),
\(\dot x=\kappa_1x^2-\kappa_2xy+\kappa_3y\), \(\dot y=2\kappa_1x^2-\kappa_2xy\) admits a limit cycle.''
\end{openquote}

The example was settled in~\autocite{Result2} with a Dulac function. Here we settle the general question.

\begin{theorem}\label{thm:A}
Let \(\mathcal N\) be a three-reaction mass-action network on two species whose sources are \(2X, X+Y, Y\) (case~9)
or \(2X, X+Y, 0\) (case~10), with arbitrary target complexes. If \(\mathcal N\) does not satisfy the Andronov--Hopf
condition of~\autocite[Theorem~9]{BBH2024}, then for every \(\kappa\in\R^3_{>0}\) its mass-action system has no
periodic orbit in \(\R^2_{\ge0}\). In particular it has no limit cycle.
\end{theorem}

\begin{corollary}\label{cor:B}
For a network of case~9 or~10, consider: (a) a periodic orbit exists for some \(\kappa\); (b) a limit cycle exists
for some \(\kappa\); (c) BBH's Andronov--Hopf condition holds. Then (a)\(\Rightarrow\)(c) by \zcref{thm:A}, and
(b)\(\Rightarrow\)(a) trivially. BBH state, from a computation of the first focal value, that the bifurcation is
supercritical, which gives (c)\(\Rightarrow\)(b). Granting that computation, which we did not re-check, the three
statements are equivalent.
\end{corollary}

\section{Setting}

Following~\autocite[(13)--(18)]{BBH2024}, reaction \(i\) is \(a_iX+b_iY\to(a_i+c_i)X+(b_i+d_i)Y\), with integer
\((c_i,d_i)\neq0\) and nonnegative target. The rate laws are \(\kappa_i m_i\), where
\(m=(x^2,xy,y)\) in case~9 and \(m=(x^2,xy,1)\) in case~10. The admissible ranges are \(c_1\ge-2\), \(d_1\ge0\),
\(c_2,d_2\ge-1\), \(c_3\ge0\), and \(d_3\ge-1\) (case~9) or \(d_3\ge0\) (case~10). Let \(u=c\times d\) and
\(A:=c_1u_1-c_3u_3\). This \(A\) is a number, not to be confused with the function \(A(x)\) of §3.

BBH's condition is the following. In case~9: either \(c_1>0\), \(c_2=-1\), \(c_3>0\), \(d_1>0\), \(d_2\ge-1\),
\(d_3\ge-1\) and \(\tfrac12(d_3/c_3+d_1/c_1)<d_2/c_2<d_1/c_1\); or \(c_1>0\), \(c_2=-1\), \(c_3=0\), \(d_1>0\),
\(d_2\ge-1\), \(d_3=-1\) and \(d_2/c_2<d_1/c_1\). In case~10: \(c_1>0\), \(c_2=-1\), \(c_3>0\), \(d_1>0\),
\(d_2=-1\), \(d_3\ge0\) and \(\tfrac12(d_3/c_3+d_1/c_1)<d_2/c_2<d_1/c_1\).

We use the following facts.

\emph{(F1)} \(\R^2_{>0}\) is forward invariant: a reaction with \(c_i<0\) has \(a_i\ge1\), and similarly in \(y\). The
boundary of the quadrant contains no Jordan curve. Hence every periodic orbit in \(\R^2_{\ge0}\) lies in
\(\R^2_{>0}\).

\emph{(F2)} If \(\operatorname{rank}\Gamma=2\), a positive equilibrium exists only if
\(u\in\R^3_{>0}\cup\R^3_{<0}\), and then it is unique: \(\kappa_im_i(\bar x,\bar y)=\mu u_i\).

\emph{(F3)} Put \(x=\bar xX\), \(y=\bar yY\) and \(t=(\bar x/\mu)T\). The system becomes
\(X'=\sum_ic_iu_im_i(X,Y)\), \(Y'=\rho\sum_id_iu_im_i(X,Y)\), with \(\rho=\bar x/\bar y\in(0,\infty)\) and
equilibrium \((1,1)\). There,
\(\det J=\rho u_1u_2u_3\) (case~9) or \(2\rho u_1u_2u_3\) (case~10), and
\(\Trop J=A-\rho d_1u_1\) (case~9) or \(A+\rho d_2u_2\) (case~10).

\emph{(F4)} We use Poincar\'e index theory, the Jordan curve theorem and the Bendixson--Dulac criterion; see
e.g.~\autocite[Ch.~3]{Perko2001}.

\section{Reduction to three families}

\begin{lemma}\label{lem:classes}
There is no periodic orbit in \(\R^2_{\ge0}\), for any \(\kappa\), in each of the following cases:
(i) \(\operatorname{rank}\Gamma\le1\);
(ii) \(\operatorname{rank}\Gamma=2\) and \(u\notin\R^3_{>0}\cup\R^3_{<0}\);
(iii) \(u\in\R^3_{<0}\);
(iv) \(u\in\R^3_{>0}\) and \(c_1\le0\);
(v) case~9, \(u\in\R^3_{>0}\), \(c_1>0\) and \(d_1=0\).
\end{lemma}

\begin{proof}
(i) Trajectories lie on lines.

(ii) \(u^\perp=\operatorname{im}\Gamma^\top\) contains some \(w\ge0\), \(w\neq0\). If \(u_i=0\) take
\(w=e_i\); if \(u_i>0>u_j\) take \(w=|u_j|e_i+u_ie_j\). Then \(\lambda\cdot(x,y)\) with \(\Gamma^\top\lambda=w\) is
strictly increasing on \(\R^2_{>0}\).

(iii) The unique positive equilibrium is a saddle, of index \(-1\), and it is the only equilibrium inside any
periodic orbit in the convex set \(\R^2_{>0}\). This contradicts index theory.

(iv) In case~9, \(\Divop(F/(xy))=(-d_1\kappa_1x^3+c_1\kappa_1x^2y-c_3\kappa_3y^2)/(x^2y^2)\). In case~10 it is
\((-d_1\kappa_1x^3+c_1\kappa_1x^2y-c_3\kappa_3y-d_3\kappa_3x)/(x^2y^2)\). Both are negative on \(\R^2_{>0}\): if
\(c_1<0\) the middle term is negative, and if \(c_1=0\) then \(d_1\ge1\). This is~\autocite[Lemma~4]{BBH2024}.

(v) Here \(u>0\) forces \(d_2\ge1\), \(d_3=-1\), \(c_2=-1\) and \(c_3=0\), and then \(\Divop(F/(xy))=c_1\kappa_1/y>0\)
(cf.~\autocite[§4.2]{BBH2024}).
\end{proof}

\begin{lemma}\label{lem:hopf}
If \(u\in\R^3_{>0}\), \(c_1>0\), \(d_1>0\) and \(A>0\), then BBH's condition holds.
\end{lemma}

\begin{proof}
Case~9. Suppose \(c_2\ge0\). Then \(u_3>0\) gives \(d_2\ge1\), and \(u_1>0\) gives \(c_2,d_3>0\). If \(c_3=0\), then
\(u_2<0\). If \(c_3>0\), then \(d_1<c_1d_2/c_2\) and \(d_3>c_3d_2/c_2\), so \(u_2=c_3d_1-c_1d_3<0\). Both contradict
\(u>0\), so \(c_2=-1\). Then \(u_3>0\) reads \(d_2/c_2<d_1/c_1\). If \(c_3>0\), \(A>0\) reads
\(\tfrac12(d_3/c_3+d_1/c_1)<d_2/c_2\). If \(c_3=0\), \(u_2>0\) gives \(d_3=-1\).

Case~10. Some \(d_i<0\) is needed, so \(d_2=-1\). Then \(u_3>0\) forces \(c_2=-1\), and \(u_1>0\) gives
\(c_3>d_3\ge0\). Conclude as in case~9.
\end{proof}

\begin{lemma}\label{lem:families}
If \(u\in\R^3_{>0}\), \(c_1>0\), \(d_1>0\) and \(A\le0\), then \(c_1,c_3\ge1\), \(c_2=-1\), and the network is in
exactly one of:
\begin{description}
\item[9a] \(2X\to(2+c_1)X+d_1Y\), \(X+Y\to0\), \(Y\to c_3X+(1+d_3)Y\), with \(d_1>c_1\), \(-1\le d_3<c_3\) and
  \(d_1/c_1+d_3/c_3\ge2\);
\item[9b] \(2X\to(2+c_1)X+d_1Y\), \(X+Y\to Y\), \(Y\to c_3X\), with \(c_3d_1\ge c_1\);
\item[10a] \(2X\to(2+c_1)X+d_1Y\), \(X+Y\to0\), \(0\to c_3X+d_3Y\), with \(d_1>c_1\), \(0\le d_3<c_3\) and
  \(d_1/c_1+d_3/c_3\ge2\).
\end{description}
In each, \(k:=c_3u_3/(c_1u_1)\ge1\), since \(A=c_1u_1(1-k)\).
\end{lemma}

\begin{proof}
\(A\le0<c_1u_1\) gives \(c_3\ge1\), and \(c_2=-1\) as in \zcref{lem:hopf}.

In case~9, \(u_1=-d_3-c_3d_2>0\) excludes \(d_2\ge1\). If \(d_2=-1\), we get 9a. If \(d_2=0\), then \(d_3=-1\) and
we get 9b.

In case~10, \(d_2=-1\) and we get 10a.
\end{proof}

\begin{proof}[of \zcref{thm:A}]
In case~10, \(u>0\) and \(c_1>0\) already force \(d_1>0\). A network failing BBH's condition is therefore in
\zcref{lem:classes}, or in a family of \zcref{lem:families} by \zcref{lem:hopf}. The families are treated by
\zcref{thm:families}.
\end{proof}

\section{The families: a Li\'enard argument}

\begin{theorem}\label{thm:families}
Networks of 9a, 9b and 10a have no periodic orbit in \(\R^2_{\ge0}\), for every \(\kappa\).
\end{theorem}

Put \(\alpha_1=c_1u_1\), \(\alpha_2=u_2\), \(\alpha_3=c_3u_3\), \(\beta_1=d_1u_1\), \(\beta_2=-d_2u_2\ge0\) and
\(\beta_3=d_3u_3\). Then \(\alpha_2=\alpha_1+\alpha_3\), \(\beta_2=\beta_1+\beta_3\) and \(k=\alpha_3/\alpha_1\ge1\).
By (F3) the system reads
\[
\dot x=A(x)+B(x)y,\qquad \dot y=\rho\bigl(C(x)+D(x)y\bigr),
\]
with
\begin{itemize}
\item case~9: \(A=\alpha_1x^2\), \(B=\alpha_3-\alpha_2x\), \(C=\beta_1x^2\), \(D=\beta_3-\beta_2x\);
\item case~10: \(A=\alpha_1x^2+\alpha_3\), \(B=-\alpha_2x\), \(C=\beta_1x^2+\beta_3\), \(D=-\alpha_2x\).
\end{itemize}

One checks:
\begin{enumerate}
\item \(BC-AD=u_1u_2u_3\,x^2(1-x)\) in case~9, and \(u_1u_2u_3\,x(1-x^2)\) in case~10.
\item On \(I=(x_0,\infty)\), with \(x_0=\alpha_3/\alpha_2\in(0,1)\) (case~9) or \(x_0=0\) (case~10), we have \(B<0\)
  and \(D/B>0\). In case~9, \(D(x_0)=-u_1u_3\) and \(D'=-\beta_2\le0\).
\item The \(x\)-nullcline \(N=-A/B\) satisfies \(N(1)=1\) and \(N'(1)\le0\).
\end{enumerate}
Periodic orbits lie in \(I\times\R\). In case~9, at a minimum \(m\le x_0\) of \(x\) we would have
\(\dot x\ge\alpha_1m^2>0\).

\begin{lemma}\label{lem:lienard}
Put \(f=N'-\rho D/B\), \(F=\int_1^xf\), \(g=-\rho(BC-AD)/B^2\), \(G=\int_1^xg\) and \(Y=y-N(x)+F(x)\). Then
\(\dot x/B=Y-F(x)\) and \(\dot Y/B=-g(x)\). Since \(B<0\) on \(I\), periodic orbits in \(I\times\R\) correspond to
periodic orbits of \(x'=Y-F(x)\), \(Y'=-g(x)\). Here \(F=N-1-\rho\psi\) with \(\psi=\int_1^xD/B\) increasing,
\(f(1)<0\), \((x-1)g(x)>0\) on \(I\setminus\{1\}\), and \(G=\rho\Gt\) with \(\Gt\) independent of~\(\rho\).
\end{lemma}

\begin{lemma}\label{lem:comparison}
Let \(I\ni\bar x\) be an open interval, \(F\in C^1(I)\) with \(F(\bar x)=0\) and \(f=F'\), and \(g\in C^1(I)\) with
\((x-\bar x)g(x)>0\) for \(x\neq\bar x\). Let \(G=\int_{\bar x}^xg\). Assume \(f(\bar x)<0\), and \(F(a)>F(b)\)
whenever \(a<\bar x<b\) in \(I\) and \(G(a)=G(b)\). Then \(x'=Y-F(x)\), \(Y'=-g(x)\) has no periodic orbit in
\(I\times\R\).
\end{lemma}

\begin{proof}
Suppose a periodic orbit \(\Gamma\) exists.

\emph{\(\Gamma\) meets \(\{x=\bar x\}\) exactly twice.} Its crossings of the line are transversal, finite in number
and at least one; otherwise \(Y\) would be monotone. They split \(\Gamma\) into right and left arcs. Each arc is a
graph \(x=X(Y)\) over an interval containing~\(0\), because \(Y\) is strictly monotone on it. Disjoint arcs on the
same side are strictly nested, by the intermediate value theorem. If there were two right arcs, the left arc leaving
the bottom of the outermost right arc and the left arc arriving at its top would be distinct and not nested, which
is impossible. So \(\Gamma=R\cup L\), both arcs being graphs over the same interval \([Y_-,Y_+]\), with
\(Y_-<0<Y_+\).

\emph{Comparison.} Let \(Z_{R,L}=G(X_{R,L}(Y))\). Then \(Z'=F(X)-Y\), so \(\Delta:=Z_L-Z_R\) satisfies
\(\Delta'=F(X_L)-F(X_R)\) and \(\Delta(Y_\pm)=0\). Since \(f(\bar x)<0\), we have \(F(X_L)>0>F(X_R)\) near both ends.
Hence \(\Delta>0\) near \(Y_-\) and \(\Delta<0\) near \(Y_+\). At the first zero \(Y^*\) of \(\Delta\) after
\(Y_-\), we have \(\Delta'(Y^*)\le0\). But there \(G(X_L)=G(X_R)\), so the hypothesis gives
\(\Delta'(Y^*)=F(X_L)-F(X_R)>0\), a contradiction.
\end{proof}

Since \(F(a)-F(b)=[N(a)-N(b)]+\rho[\psi(b)-\psi(a)]\) and the second bracket is positive, the hypothesis of
\zcref{lem:comparison} holds for \emph{all}~\(\rho\) once
\begin{equation}\label{eq:star}
N(a)\ge N(b)\quad\text{whenever } a<1<b \text{ in } I \text{ and } \Gt(a)=\Gt(b).
\end{equation}

\begin{lemma}\label{lem:case9}
\eqref{eq:star} holds in 9a and 9b.
\end{lemma}

\begin{proof}
With \(x=(\eta+k)/(1+k)\), \(\eta\in(0,\infty)\), we get \(N=(\eta+k)^2/((1+k)^2\eta)\) and \(\Gt=K\,H_k(\eta)\)
with \(K>0\), where
\[
H_k(\eta)=\int_1^\eta\frac{(\sigma+k)^2(\sigma-1)}{\sigma^2}\,d\sigma
=\frac{\eta^2}{2}+(2k-1)\eta+(k^2-2k)\log\eta+\frac{k^2}{\eta}-\Bigl(k^2+2k-\frac12\Bigr).
\]
Let \(P(a)=H_k(k^2/a)-H_k(a)\). Then \(P'(a)=-Q(a)/a^3\), with
\(Q(a)/a^2=8k(k-1)+(6k-2)s+s^2>0\), where \(s=(a-k)^2/a>0\) for \(0<a<1\le k\). Also \(P(1)=H_k(k^2)\ge0\). So
\(P>0\) on \((0,1)\).

Hence, if \(H_k(\eta_a)=H_k(\eta_b)\) with \(\eta_a<1<\eta_b\), then \(\eta_b<k^2/\eta_a\), i.e.
\(\eta_a\eta_b<k^2\). Therefore
\[
N(a)-N(b)=\frac{(\eta_a-\eta_b)\bigl(1-k^2/(\eta_a\eta_b)\bigr)}{(1+k)^2}>0.
\]
\end{proof}

\begin{lemma}\label{lem:case10}
\eqref{eq:star} holds in 10a.
\end{lemma}

\begin{proof}
Here \(g/\rho=(u_1u_3/u_2)\,h'\) with \(h=x^2/2-\log x\). The function \(t\mapsto h(1/t)-h(t)\) vanishes at~\(1\)
and has derivative \(-(1-t^2)^2/t^3\). So pairs satisfy \(ab<1\le k\), and
\(N(a)-N(b)=(a-b)(\alpha_1-\alpha_3/(ab))/\alpha_2>0\).
\end{proof}

\begin{proof}[of \zcref{thm:families}]
Combine (F1), the confinement to \(I\times\R\), and Lemmas~\ref{lem:lienard}--\ref{lem:case10}.
\end{proof}

\begin{remark}
In case~10 the computation of \zcref{lem:case10} works for every \(k\). For \emph{every} case-10 network with
\(u>0\) and \(c_1>0\), Hopf networks included,
\(F(a)-F(b)=(b-a)[\rho-(\alpha_1-\alpha_3/(ab))/\alpha_2]\), and \(\Trop J<0\iff\rho>N'(1)\). Hence there is no
periodic orbit whenever \(\Trop J<0\). For case-9 Hopf networks we observed the analogous sharpness numerically, but
we have not proved it.

In the subfamily of~\autocite{Result2} (9a with \(d_3=0\)), \(k\ge1\) is \(d_1\ge2c_1\).
\end{remark}

\section{Computer verification}

\begin{itemize}
\item \texttt{scripts/certify\_lienard.py} (sympy, exact; 64 checks) verifies:
  \begin{itemize}
  \item the normalisation, the equilibrium formulas, \(\det J\) and \(\Trop J\);
  \item the Li\'enard identity for general \(A,B,C,D\);
  \item for the symbolic integer data of each family: \(u>0\), (1)--(3), \(f(1)<0\), and
    \(A=\alpha_1-\alpha_3\) equal to BBH's inequality;
  \item the \(\eta\)-substitution, the formula for \(H_k\) and for \(P'\), the decomposition of \(Q\), the
    case-10 identities;
  \item the Dulac divergences.
  \end{itemize}
\item \texttt{scripts/classify\_networks.py} enumerates all case-9 and case-10 networks whose target coefficients are
  at most 20 (\(1.7\cdot10^8\) networks). It confirms that each one falls in exactly one class of
  Lemmas~\ref{lem:classes}--\ref{lem:families} or is a BBH Hopf network, and that BBH's condition is equivalent to
  \(u>0\), \(c_1>0\), \(d_1>0\), \(A>0\).
\item \texttt{scripts/numerical\_sanity.py} (floating point, not part of the proof) finds no cycle in 128 backward
  integrations of random family members, finds the cycles of Hopf networks past the bifurcation, and checks the
  hypothesis of \zcref{lem:comparison} by quadrature.
\end{itemize}

\paperreferences
\end{document}
